% Надпреварата за собственост върху кадъра на корутината в await_suspend.
%
% Времева диаграма с две ленти. Времето тече надолу, а хоризонталното разстояние
% е само разделение между нишките, не мащаб. Двете ленти са центрирани на 3.3 cm
% и 10.9 cm, кутиите са широки 4.3 cm, така че пълната ширина излиза около
% 13.2 cm: побира се в 14.5 cm без мащабиране.
%
% Черно бяло: опасните стъпки се различават по дебела рамка, стъпките на
% поправката по сив пълнеж, а прозорецът на надпреварата по сивото поле с
% пунктирана граница отгоре. Никъде няма цвят.
%
% Двата панела са нарочно с еднаква геометрия, за да се чете поправката като
% размяна на две стъпки, а не като друг алгоритъм.
%
% Забележка към редактора: в възел с align=center преминаването на нов ред \\
% не може да е вътре в група. Затова основният шрифт на кутиите е този на
% пояснението, а редът с кода е обвит в собствена група с по-голям кегел.
\begin{figure}[htbp]
\centering
\begin{tikzpicture}[
  ev/.style={box, fill=white, minimum width=4.3cm, minimum height=0.78cm,
             font=\fontsize{7}{8.5}\selectfont},
  hz/.style={ev, very thick},                 % достъп до кадър, който може да е изчезнал
  key/.style={ev, fill=gray!25, very thick},  % двете стъпки на поправката
  hdr/.style={align=center, font=\fontsize{8}{9.5}\selectfont},
  ttl/.style={align=center, font=\fontsize{9}{11}\bfseries\selectfont},
  gap/.style={align=center, text width=3.0cm, font=\fontsize{7.5}{9}\selectfont},
  edge label/.style={font=\fontsize{7}{8.5}\selectfont},
  lane/.style={densely dotted, semithick, -{Stealth[length=2mm]}},
]

% ===== Панел 1: наредбата, която въвежда надпреварата ======================
\node[ttl] at (7.1,0.45) {Наредба, която въвежда надпревара};

\node[hdr] at (3.3,-0.10)  {\textbf{Нишка A}, подаваща\\ вътре в \code{await\_suspend}};
\node[hdr] at (10.9,-0.10) {\textbf{Нишка B}, подновяваща\\ обработва завършването};

\node[ev] (a1) at (3.3,-1.10)
     {{\fontsize{8}{9.5}\selectfont\code{op.continuation = h;}}\\ продължението е записано};
\node[ev] (a2) at (3.3,-2.35)
     {{\fontsize{8}{9.5}\selectfont\code{submit()}}\\ подаването се извършва};
\node[hz] (a3) at (3.3,-3.60)
     {{\fontsize{8}{9.5}\selectfont\code{submitted = ...;}}\\ запис в член на кадъра};
\node[hz] (a4) at (3.3,-4.85)
     {{\fontsize{8}{9.5}\selectfont\code{return submitted;}}\\ четене на член на кадъра};

% Пълнежът е в тона на полето, за да прекъсне пунктираната лента под текста.
\node[align=center, font=\fontsize{7.5}{9}\selectfont, fill=gray!25,
      inner sep=2pt] (an) at (3.3,-5.95)
     {Двата достъпа по-горе са в кадър,\\ който може вече да не съществува.};

\node[ev] (b1) at (10.9,-3.40)
     {{\fontsize{8}{9.5}\selectfont завършването пристига}\\ възможно е още тук};
\node[ev] (b2) at (10.9,-4.60)
     {{\fontsize{8}{9.5}\selectfont\code{h.resume()}}\\ корутината се подновява};
\node[ev] (b3) at (10.9,-5.80)
     {{\fontsize{8}{9.5}\selectfont кадърът се унищожава}\\ паметта му е освободена};

\draw[flow] (a1) -- (a2);
\draw[flow] (a2) -- (a3);
\draw[flow] (a3) -- (a4);
\draw[flow] (b1) -- (b2);
\draw[flow] (b2) -- (b3);

% Подаването е точката, в която операцията става годна за завършване. Оттук
% нататък двете ленти вървят едновременно и наредбата между тях не е гарантирана.
% Стрелката е с прав ъгъл, а не диагонална, за да не мине през кутията вдясно.
\draw[thin flow] (a2.east) -| (b1.north);
\node[edge label, above] at (8.2,-2.37) {годна за завършване};

% Обратната пунктирана стрелка казва защо горните два достъпа са невалидни.
\draw[thin flow] (b3.west) -- (an.east);

\node[gap] at (7.1,-4.70)
     {\textbf{Прозорец на надпреварата}\\[2pt]
      Оттук нататък кадърът е на подновяващата нишка.};

% ===== Панел 2: наредбата, която издържа ==================================
\node[ttl] at (7.1,-7.25) {Наредба, която издържа};

\node[hdr] at (3.3,-7.85)  {\textbf{Нишка A}, подаваща};
\node[hdr] at (10.9,-7.85) {\textbf{Нишка B}, подновяваща};

\node[ev] (c1) at (3.3,-8.65)
     {{\fontsize{8}{9.5}\selectfont\code{op.continuation = h;}}\\ продължението е записано};
\node[key] (c2) at (3.3,-9.95)
     {{\fontsize{8}{9.5}\selectfont\code{submitted = true;}}\\ публикувано, докато\\ кадърът още е наш};
\node[ev] (c3) at (3.3,-11.40)
     {{\fontsize{8}{9.5}\selectfont\code{if (!submit())}}\\ подаването е последното\\ докосване на кадъра};
\node[key] (c4) at (3.3,-12.80)
     {{\fontsize{8}{9.5}\selectfont\code{return true;}}\\ литерал, а не член};

\node[ev] (d1) at (10.9,-12.80)
     {{\fontsize{8}{9.5}\selectfont завършването може да поднови}\\ корутината и кадърът да изчезне};

\draw[flow] (c1) -- (c2);
\draw[flow] (c2) -- (c3);
\draw[flow] (c3) -- (c4);

% Стрелката е с прав ъгъл, а не диагонална, за да не пресече етикета в средата.
\draw[thin flow] (c3.east) -| (d1.north);
\node[edge label, above] at (8.2,-11.42) {годна за завършване};

\node[gap] at (7.1,-12.80)
     {Тук няма достъп до кадъра, значи няма и какво да се надпреварва.};

% ===== Фон: прозорците и лентите =========================================
\begin{scope}[on background layer]
  \fill[gray!25] (0.85,-2.90) rectangle (13.35,-6.40);
  \fill[gray!10] (0.85,-12.10) rectangle (13.35,-13.55);
  % Пунктираната граница отбелязва момента, а сивото полето след него.
  \draw[thick, densely dashed] (0.85,-2.90) -- (13.35,-2.90);
  \draw[thick, densely dashed] (0.85,-12.10) -- (13.35,-12.10);

  % Лентите спират вътре в полето, за да не увиснат стрелки под него.
  \draw[lane] (3.3,-0.60)  -- (3.3,-6.30);
  \draw[lane] (10.9,-0.60) -- (10.9,-6.30);
  \draw[lane] (3.3,-8.15)  -- (3.3,-13.45);
  \draw[lane] (10.9,-8.15) -- (10.9,-13.45);
\end{scope}

\draw[-{Stealth[length=2mm]}, semithick] (0.42,-1.15) -- (0.42,-2.85);
\node[rotate=90, font=\fontsize{7.5}{9}\selectfont, anchor=south] at (0.32,-2.00) {време};

\end{tikzpicture}
\caption[Надпреварата за кадъра на корутината в \code{await\_suspend}]%
{Надпреварата за собственост върху кадъра на корутината, раздел~\ref{sec:awaiter}.
Горе: подаването прави операцията годна за завършване, затова всеки следващ достъп
до член на кадъра от подаващата нишка е достъп до евентуално освободена памет.
Долу: поправката публикува стойността преди подаването и връща литерал. Пътят на
неуспеха, който записва \code{false} повторно, е безопасен по устройство, защото
щом нищо не е подадено, нищо не може да подновява.}
\label{fig:awaiter_race}
\end{figure}
